\documentclass[preview]{standalone}

\usepackage{amsmath}
\usepackage{amssymb}
\usepackage{stellar}
\usepackage{definitions}

\begin{document}

\id{ring-homomorphisms}
\genpage

\section{Ring homomorphisms}

\begin{snippetdefinition}{ring-homomorphism-definition}{Ring homomorphism}
    Let \(A\) and \(B\) be \ring[rings]. A \function
    \(\varphi\colon A\fromto B\) is a \emph{ring homomorphism} if, for all
    \(a_1,a_2\in A\),
    \[
        \varphi(a_1+a_2)=\varphi(a_1)+\varphi(a_2),
        \qquad
        \varphi(a_1a_2)=\varphi(a_1)\varphi(a_2),
    \]
    and
    \[
        \varphi(1_A)=1_B.
    \]
\end{snippetdefinition}

\begin{snippetdefinition}{ring-kernel-definition}{Kernel of a ring homomorphism}
    Let \(\varphi\colon A\fromto B\) be a \ringhomomorphism. Its
    \emph{kernel} is the \set
    \[
        \operatorname{ker}(\varphi)
        =\{a\in A\suchthat\varphi(a)=0_B\}.
    \]
\end{snippetdefinition}

\begin{snippetdefinition}{ring-isomorphism-definition}{Ring isomorphism}
    Let \(A\) and \(B\) be \ring[rings]. A \emph{ring isomorphism} from
    \(A\) to \(B\) is a \bijective \ringhomomorphism
    \(\varphi\colon A\fromto B\). If such a \function exists, one writes
    \[
        A\cong B.
    \]
\end{snippetdefinition}

\begin{snippetproposition}{ring-homomorphism-properties}{Elementary properties of ring homomorphisms}
    Let \(\varphi\colon A\fromto B\) be a \ringhomomorphism.
    \begin{enumerate}
        \item If \(a\in A^\times\), then
        \(\varphi(a)\in B^\times\) and
        \(\varphi(a^{-1})=\varphi(a)^{-1}\).
        \item If \(S\subringleq A\), then
        \(\varphi(S)\subringleq B\).
        \item If \(R\subringleq B\), then
        \[
            \varphi^{-1}(R)
            =\{a\in A\suchthat\varphi(a)\in R\}
            \subringleq A.
        \]
        \item \(\ringker\varphi\idealin A\).
        \item The \function \(\varphi\) is \injective \ifandonlyif
        \[
            \ringker\varphi=\{0_A\}.
        \]
    \end{enumerate}
\end{snippetproposition}

\begin{snippetproof}{ring-homomorphism-properties-proof}{ring-homomorphism-properties}{Elementary properties of ring homomorphisms}
    If \(ab=ba=1_A\), then
    \[
        \varphi(a)\varphi(b)=\varphi(b)\varphi(a)=1_B,
    \]
    proving the first assertion.

    Let \(S\subringleq A\). Since \(1_B=\varphi(1_A)\in\varphi(S)\), and
    \[
        \varphi(s_1)-\varphi(s_2)=\varphi(s_1-s_2),
        \qquad
        \varphi(s_1)\varphi(s_2)=\varphi(s_1s_2),
    \]
    the image \(\varphi(S)\) is a \subring. The same identities show that
    the preimage of any \subring of \(B\) is a \subring of \(A\).

    The kernel contains \(0_A\) and is closed under differences. If
    \(x\in\ringker\varphi\) and \(a\in A\), then
    \[
        \varphi(ax)=\varphi(a)0_B=0_B,
        \qquad
        \varphi(xa)=0_B\varphi(a)=0_B.
    \]
    Hence \(\ringker\varphi\idealin A\).

    If \(\varphi\) is \injective, then
    \(\varphi(a)=0_B=\varphi(0_A)\) implies \(a=0_A\). Conversely, if
    \(\ringker\varphi=\{0_A\}\) and
    \(\varphi(a_1)=\varphi(a_2)\), then
    \(a_1-a_2\in\ringker\varphi\), so \(a_1=a_2\).
\end{snippetproof}

\begin{snippetproposition}{ring-homomorphism-composition}{Composition of ring homomorphisms}
    Let
    \[
        A\xrightarrow{\varphi}B\xrightarrow{\psi}C
    \]
    be \ringhomomorphism[ring homomorphisms]. Then
    \[
        \psi\circ\varphi\colon A\fromto C
    \]
    is a \ringhomomorphism.
\end{snippetproposition}

\begin{snippetproof}{ring-homomorphism-composition-proof}{ring-homomorphism-composition}{Composition of ring homomorphisms}
    Both addition and multiplication are preserved because each of
    \(\varphi\) and \(\psi\) preserves them. Moreover,
    \[
        (\psi\circ\varphi)(1_A)
        =\psi(\varphi(1_A))
        =\psi(1_B)
        =1_C.
    \]
\end{snippetproof}

\end{document}
